\subsection{Solution Methodology}
The branch and bound method is used to solve integer linear programming problems (ILP) by iteratively restricting the feasible region of a linear programming relaxation of the problem. The method involves branching on decision variables, creating subproblems, and using bounds to eliminate subproblems that cannot yield better solutions than the best known solution. The process continues until all subproblems have been explored or eliminated, resulting in the optimal integer solution.

As an example, cosider the following ILP problem from the 2022 exam:

\begin{align*}
    \max. \quad & 55x_1 + 53x_2 + 81x_3 \\
    \text{s.t.} \quad & 8x_1 + 2x_2 + 4x_3 \le 50 \\
    & 1x_1 + 8x_2 + 7x_3 \le 60 \\
    & 4x_1 + 5x_2 + 4x_3 \le 50 \\
    & x_1, x_2, x_3 \ge 0 \text{ og heltallig}
\end{align*}

Here we get the branching tree as follows:

\begin{tikzpicture}[
    level distance=1.5cm,
    sibling distance=4cm,
    every node/.style={circle, draw, inner sep=2pt, minimum size=6mm, font=\small},
    label distance=1mm,
    edge from parent/.style={draw, -},
    edge from parent path={(\tikzparentnode) -- (\tikzchildnode)}
]

% Adjust sibling distances for deeper levels
\tikzset{
    level 1/.style={sibling distance=6cm},
    level 2/.style={sibling distance=3cm},
    level 3/.style={sibling distance=2cm}
}

% Root Node P0
\node[label=above right:{\begin{tabular}{l}\scriptsize 786.2 \\ \scriptsize (2.1, 0, 8.3)\end{tabular}}] {$P_0$}
    % Left child P1
    child {
        node[label=above left:{\begin{tabular}{r}\scriptsize 781.1 \\ \scriptsize (2, 0, 8.3)\end{tabular}}] {$P_1$}
        % P1 Left child P2
        child {
            node[label=above left:{\begin{tabular}{r}\scriptsize 771.3 \\ \scriptsize (2, 0.3, 8)\end{tabular}}] {$P_2$}
            % P2 Left child P3
            child {
                node[label=left:{\begin{tabular}{r}\scriptsize 758* \\ \scriptsize (2, 0, 8)\end{tabular}}] {$P_3$}
                edge from parent node[draw=none, left, font=\scriptsize, xshift=-2mm] {$x_2 \le 0$}
            }
            % P2 Right child P4
            child {
                node[label=right:{\begin{tabular}{l}\scriptsize 741.6 \\ \scriptsize (2, 1, 7.1)\end{tabular}}] {$P_4$}
                edge from parent node[draw=none, right, font=\scriptsize, xshift=2mm] {$x_2 \ge 1$}
            }
            edge from parent node[draw=none, left, font=\scriptsize, xshift=-2mm] {$x_3 \le 8$}
        }
        % P1 Right child P5
        child {
            node[label=right:{\begin{tabular}{l}\scriptsize Ikke- \\ \scriptsize brugbar\end{tabular}}] {$P_5$}
            edge from parent node[draw=none, right, font=\scriptsize, xshift=2mm] {$x_3 \ge 9$}
        }
        edge from parent node[draw=none, above left, font=\scriptsize] {$x_1 \le 2$}
    }
    % Right child P6
    child {
        node[label=right:{\begin{tabular}{l}\scriptsize 723.4 \\ \scriptsize (3, 2.6, 5.2)\end{tabular}}] {$P_6$}
        edge from parent node[draw=none, above right, font=\scriptsize] {$x_1 \ge 3$}
    };

\end{tikzpicture}

In $P_0$, we branch on $x_1$ and create a subproblem $P_1$ with the constraint $x_1 \le 2$, as we adopt a depth first search strategy, we continue by branching on $x_3$ in $P_1$ and create $P_2$ with the constraint $x_3 \le 8$. In $P_2$, we branch on $x_2$ and create $P_3$ with the constraint $x_2 \le 0$. We find that the optimal solution to $P_3$ is integral, so we can update our incumbent solution to 758. We then backtrack to $P_2$ and branch on $x_2 \ge 1$ and create $P_4$. The optimal solution to $P_4$ is 741.6, which is less than our incumbent solution, hence we can prune this branch. We backtrack to $P_1$ and branch on $x_3 \ge 9$ and create $P_5$. The optimal solution to $P_5$ is not feasible, so we can fathom $P_5$. We backtrack to the root node and branch on $x_1 \ge 3$ and create $P_6$. The optimal solution to $P_6$ is 723.4, which is less than our incumbent solution, so we can fathom $P_6$. Since all nodes have been explored or fathomed, we conclude that the optimal integer solution to the original problem is 758.

